\bookchapter{Chapitre 47 : Aide-mémoire et ressources}

\chapterintro{Ce chapitre est un aide-mémoire qui se consulte de façon autonome. Quand un principe t'échappe ou qu'il te faut décider vite, ouvre-le directement. Traite-le comme une « fiche de poche » : les cinq grands principes en un clin d'œil, l'arbre de décision des outils, la FAQ, les ressources d'apprentissage — il y a tout.}

\booksection{Les cartes mémo des principes de pensée}

\begin{center}
\begin{tabularx}{\linewidth}{|X|X|X|}
\hline
\textbf{Principe} & \textbf{En une phrase} & \textbf{Points clés} \tabularnewline
\hline
Think Before Coding & Réfléchir avant de parler & Réponds aux trois questions : quoi faire ? à quoi ressemble le résultat ? quelles contraintes ? \tabularnewline
Simplicity First & Résoudre avec un minimum de code & Une seule fonctionnalité à la fois, couper la complexité inutile \tabularnewline
Surgical Changes & Modifier avec précision, sans déborder & Délimite clairement le périmètre de modification ; refuse le « j'en profite pour retoucher » de l'IA \tabularnewline
Goal-Driven Execution & Dire « quoi faire », pas « comment faire » & Donner l'objectif, pas les étapes : laisser l'IA choisir son chemin \tabularnewline
Iterate \& Verify & La boucle générer-vérifier-corriger & Vérification à trois niveaux : fonctionnel, limites, architecture \tabularnewline
\hline
\end{tabularx}
\end{center}

\begin{center}\small\color{gray}Carte mémo des cinq grands principes de pensée\end{center}

\booksection{L'arbre de décision du choix d'outils}

Quand tu dois choisir un outil de programmation IA, suis cet arbre de décision. Attention : ce n'est pas un classement du « meilleur », mais un cadre pour déterminer « lequel te convient le mieux ».

Que veux-tu faire ?

Débutant, envie d'essayer vite choisis un Builder en ligne (v0/Bolt/Replit)

Dév au quotidien, besoin d'autocomplétion choisis l'IA intégrée à l'IDE (Cursor/Trae)

Gros projets, besoin d'opérations par lots choisis un Agent en terminal (Claude Code)

Prototype rapide pour valider une idée choisis un Builder en ligne ou Trae SOLO

\begin{center}\small\color{gray}Figure 6 : Arbre de décision pour choisir un outil de programmation IA\end{center}

\booksection{La FAQ des questions courantes}

Q1 : Je n'ai jamais programmé, puis-je quand même coder avec l'IA ?

Oui. En 2026, les outils de programmation IA sont assez puissants pour que tu décrives ton besoin en langue naturelle et laisses l'IA générer le code. Deux conseils : d'abord, démarre avec des outils Builder en ligne (comme v0 ou Bolt), qui ne demandent aucune configuration d'environnement ; ensuite, apprends les concepts de base en pratiquant — pas besoin de maîtriser la syntaxe, mais il faut pouvoir « comprendre dans les grandes lignes » le code écrit par l'IA, sinon la réception ne sera pas fiable.

Q2 : Quel langage de programmation devrais-je apprendre ?

Cela dépend de ce que tu veux faire. Front-end de sites web : JavaScript/TypeScript ; traitement de données ou IA : Python ; back-end d'entreprise : Java ou Go. Mais surtout : à l'ère de la programmation IA, l'importance du « langage » décline. L'IA gère les détails de syntaxe ; ce que tu dois comprendre, c'est la « structure d'un programme » — ce qu'est une fonction, une boucle, une structure de données. Ces concepts transcendent les langages : une fois acquis, tu passes de l'un à l'autre très vite.

Q3 : Le code écrit par l'IA pose-t-il des problèmes de droits d'auteur ?

En août 2026, la question reste juridiquement controversée. En général, le code généré par l'IA n'est pas considéré comme ayant un titulaire de droits d'auteur clairement établi dans la plupart des juridictions. Mais pour un projet commercial, quelques recommandations : 1) lis les clauses de propriété intellectuelle des conditions d'utilisation de ton outil IA ; 2) relis et modifie toi-même le code critique pour y apporter ta contribution créative ; 3) évite que l'IA copie directement du code open source susceptible d'être protégé.

Q4 : L'IA va-t-elle remplacer les programmeurs ?

Pas à court terme. Ce que l'IA remplace, ce n'est pas le métier de « programmeur », c'est le travail de « codage ligne à ligne ». Comme la calculatrice n'a pas remplacé les mathématiciens, mais les a libérés des calculs pénibles à la main. Les outils de programmation IA font passer le programmeur de « l'écriture de code » à « la réflexion sur les problèmes » — ce qui, en réalité, relève son plafond de valeur. Demain, l'avantage compétitif ne sera pas de savoir « combien de lignes on écrit », mais « quels problèmes complexes on sait résoudre ».

Q5 : Comment juger si la production de l'IA est correcte ?

La vérification en trois niveaux. Premier niveau, fonctionnel : ça tourne ? Entrées et sorties sont-elles correctes ? C'est la base. Deuxième niveau, limites : que se passe-t-il avec une entrée vide, trop longue ou anormale ? Cela révèle les problèmes cachés. Troisième niveau, architecture : la structure du code est-elle saine ? Du code en dur ? Facile à faire évoluer ? Ce niveau demande un certain jugement technique. Si tu ne peux pas l'évaluer, fais au moins les deux premiers — un programme qui tourne et tient les cas limites te place déjà au-dessus de la plupart des débutants.

\booksection{Les ressources d'apprentissage recommandées}

Voici les ressources recommandées à la date d'août 2026. Attention : les liens du web évoluent ; si un lien est cassé, une recherche par son titre mène généralement à la version la plus récente.

\begin{center}
\begin{tabularx}{\linewidth}{|X|X|X|}
\hline
\textbf{Type} & \textbf{Ressource} & \textbf{Pourquoi la recommander} \tabularnewline
\hline
Sources primaires & Le CLAUDE.md d'Andrej Karpathy & La source des cinq grands principes : 70 lignes, mais chaque mot compte \tabularnewline
Sources primaires & Documentation officielle d'Anthropic (Claude Code) & Le guide des bonnes pratiques pour les agents en terminal \tabularnewline
Tutoriels & La collection de Skills assemblée par la communauté Vibe Coding & Un recueil d'astuces de terrain résumées par les praticiens de la communauté \tabularnewline
Tutoriels & La documentation officielle et la chaîne YouTube de Cursor & La référence pour l'IA intégrée à l'IDE \tabularnewline
Communautés & awesome-claude-code sur GitHub & Bonnes pratiques et plugins pour Claude Code \tabularnewline
Communautés & La communauté r/cursor de Reddit & Retours d'expérience et discussions des utilisateurs de Cursor du monde entier \tabularnewline
Bases & MDN Web Docs (developer.mozilla.org) & La documentation de référence du développement front-end \tabularnewline
Bases & Le tutoriel officiel de Python (docs.python.org) & Le point de départ officiel pour apprendre Python \tabularnewline
\hline
\end{tabularx}
\end{center}

\begin{center}\small\color{gray}La liste des ressources d'apprentissage recommandées\end{center}

\booksection{Conclusion : les outils en main, la méthode en tête}

Revenons à la conviction centrale de ce livre : les outils se renouvellent tous les trois mois, la méthodologie n'évolue que tous les trois ans. En ce mois d'août 2026, les outils de programmation IA sont devenus d'une puissance étonnante — tu peux créer un site en une phrase, générer une application full-stack à partir d'une description, mener une refactorisation complexe en une seule instruction. Mais plus les outils sont puissants, plus la méthodologie compte — parce que leur puissance amplifie chacune de tes décisions.

Si ton besoin est clair, l'IA l'implémente efficacement ; s'il est flou, l'IA fabrique efficacement la mauvaise chose. Si tu restes fidèle à la simplicité, l'IA valide vite tes idées ; si tu veux tout et tout de suite, elle entasse vite la dette technique. Si tu modifies avec précision, elle optimise ton code comme un chirurgien ; si tes instructions sont floues, elle introduit des bugs à la machette.

Voilà pourquoi ce livre consacre l'essentiel de ses pages à la méthodologie plutôt qu'aux manipulations d'outils. Non que les outils soient secondaires, mais la méthodologie est le « multiplicateur » des outils : meilleure est ta méthode, plus tes outils valent ; à méthodologie nulle, même l'outil le plus puissant ne vaut que zéro fois n'importe quoi.

Tu ne programmes plus seul : tu es le commandant, l'IA est ta légion. La qualité de ta réflexion détermine la puissance de combat de ta légion. À partir d'aujourd'hui, arme ta pensée avec la méthodologie et amplifie tes capacités avec les outils – c'est tout le secret de la programmation par IA.

\booksection{Bibliographie}

\begin{enumerate}
  \item Andrej Karpathy. (2025). Vibe Coding. Social media post introducing the concept of vibe-driven programming with AI.
  \item Andrej Karpathy. (2025). CLAUDE.md - Project guidelines for AI-assisted development. GitHub repository, 70+K stars.
  \item Anthropic. (2026). Claude Code: Terminal-native AI programming. Official documentation and best practices.
  \item Cursor Team. (2026). Cursor 3: The next generation of AI-first IDE. Release notes and documentation.
  \item OpenAI. (2026). Codex CLI: Bringing AI to the terminal. Technical documentation.
  \item Model Context Protocol (MCP) Specification. (2025-2026). Open protocol for AI-tool integration.
  \item Vibe Coding Community. (2026). Community-curated skills and best practices for AI-assisted programming.
  \item Spec-Driven Development (SDD) methodology. (2026). Emerging practices in specification-first AI programming.
\end{enumerate}
